\documentclass[12pt]{article}
\usepackage[T1]{fontenc}
\usepackage[utf8]{inputenc}
\usepackage{lmodern}
\usepackage{textcomp}
\usepackage{enumitem}
\usepackage{geometry}
\usepackage{graphicx}
\usepackage{amsmath, amssymb}
\geometry{margin=1in}
\begin{document}
\begin{center}
History - Undergraduate Level Assessment 11 \
Total Marks: 40 \
Answer all questions.
\end{center}

\section*{Multiple Choice (10 marks)}
\begin{enumerate}[label=\arabic*.]
\item Homo habilis emerged around \_\_\_\_\_\_\_\_ years ago; Homo erectus followed at around \_\_\_\_\_\_\_\_\_ years ago.
\begin{enumerate}[label=(\alph*)]
    \item 1.8 million; 1.2 million
    \item 2.5 million; 1.8 million
    \item 1.4 million; 800,000
    \item 4.5 million; 3.2 million
    \item 2 million; 1.5 million
\end{enumerate}
\item What do the amounts of rainfall in the Maya homeland, as reflected in variations in oxygen isotopes in a stalagmite in Yok Balum Cave in Belize, correspond to?
\begin{enumerate}[label=(\alph*)]
    \item severe drought and a surge in warfare
    \item increased rainfall and population migration to the Maya ceremonial centers
    \item severe drought and increased storage of surplus food from agriculture
    \item severe drought and a decline in construction at Maya ceremonial centers
    \item increased warfare and the abandonment of the entire region
\end{enumerate}
\item There is evidence that the Hohokam borrowed which practice from Mesoamerican culture?
\begin{enumerate}[label=(\alph*)]
    \item use of a calendar system based on celestial bodies
    \item practice of mummification for the dead
    \item ritual consumption of maize beer
    \item construction of underground dwellings
    \item ritual human sacrifice
\end{enumerate}
\item During the Pleistocene era, the islands of Java, Sumatra, Bali, and Borneo formed \_\_\_\_\_\_\_\_\_\_\_. The landmass connecting Australia, New Guinea, and Tasmania was called \_\_\_\_\_\_\_\_\_\_\_.
\begin{enumerate}[label=(\alph*)]
    \item Sunda; Sahul
    \item Sunda; Wallacea
    \item Gondwana; Beringia
    \item Sunda; Gondwana
    \item Beringia; Sunda
\end{enumerate}
\item Which of the following indicates a primary context?
\begin{enumerate}[label=(\alph*)]
    \item the discovery of an artifact in a museum or private collection
    \item an archaeological site found underwater
    \item a site where artifacts have been intentionally buried for preservation
    \item all of the above
    \item a site where artifacts have been moved and rearranged by humans or natural forces
\end{enumerate}
\end{enumerate}

\section*{True / False (10 marks)}
\textit{Answer True or False}
\begin{enumerate}[label=\arabic*.]
\setcounter{enumi}{5}
\item True or False: He had been replaced by Abercrombie as commander in chief after the failures of 1757. Pitt's plan called for three major offensive actions involving large numbers of regular troops, supported by the provincial militias, aimed at capturing the heartlands of New France. a different answer falling to sizable British forces.
\item True or False: Macao: The event was held in Macau on a different answer.
\item True or False: Martin Luther, who initiated the a different answer, said: "Mother Mary, like us, was born in sin of sinful parents, but the Holy Spirit covered her, sanctified and purified her so that this child was born of flesh and blood, but not with sinful flesh and blood.
\item True or False: Captain John Charles Marshall and Thomas Gilbert visited the islands in 1788. The islands were named for Marshall on Western charts, although the natives have historically named their home "jolet jen Anij" (Gifts from God). Around 1820, Russian explorer Adam Johann von Krusenstern and the French explorer Louis Isidore Duperrey named the islands after John Marshall, and drew maps of the islands.
\item True or False: Opportunistic bands of Normans successfully established a foothold in Southern Italy (the Mezzogiorno). Probably as the result of returning pilgrims' stories, the Normans entered the Mezzogiorno as warriors in 1017 at the latest. In 1018, according to Amatus of Montecassino, Norman pilgrims returning from Jerusalem called in at the port of Salerno when a Saracen attack occurred.
\end{enumerate}

\section*{Long Form Response (20 marks)}
\textit{Answer each question with a clear and complete explanation.}
\begin{enumerate}[label=\arabic*.]
\setcounter{enumi}{10}
\item Read the following passage and answer the question: Everton were founder members of the Premier League in 1992, but struggled to find the right manager. Howard Kendall had returned in 1990 but could not repeat his previous success, while his successor, Mike Walker, was statistically the least successful Everton manager to date. When former Everton player Joe Royle took over in 1994 the club's form started to improve; his first game in charge was a 2-0 victory over derby rivals Liverpool. Royle dragged Everton clear of relegation, leading the club to the FA Cup for the fifth time in its history, defeating Manchester United 1-0 in the final. Question: What year did Howard Kendall return to manage Everton?
\item Read the following passage and answer the question: The Semitic Akkadian language is first attested in proper names of the kings of Kish c. 2800 BC, preserved in later king lists. There are texts written entirely in Old Akkadian dating from c. 2500 BC. Use of Old Akkadian was at its peak during the rule of Sargon the Great (c. 2270-2215 BC), but even then most administrative tablets continued to be written in Sumerian, the language used by the scribes. Gelb and Westenholz differentiate three stages of Old Akkadian: that of the pre-Sargonic era, that of the Akkadian empire, and that of the "Neo-Sumerian Renaissance" that followed it. Akkadian and Sumerian coexisted as vernacular languages for about one thousand years, but by around 1800 BC, Sumerian was becoming more of a literary language familiar mainly only to scholars and scribes. Thorki... Question: Where is the Semetic Akkadian language first found?
\end{enumerate}

\vfill
\noindent\textit{End of Paper}
\end{document}